% Proof module for the exact common-record Fisher supremum in the shared-kernel qutrit.
% Intended for later inclusion in the Supplemental Material.

\subsection{Exact common-record Fisher supremum for the shared-kernel qutrit}
\label{supp:wp15-exact}

Use the fixed-total-excitation qutrit benchmark defined in the main supplement. Choose an orthonormal basis $\{|e_1\rangle,|e_2\rangle\}$ for the support of $\rho_0$ such that, in the ordered basis $\{|e_1\rangle,|e_2\rangle,|q\rangle\}$,
\begin{equation}
\rho_0=\operatorname{diag}(1/2,1/2,0)
\end{equation}
and
\begin{equation}
A=
\begin{pmatrix}
0 & \sqrt2/2 & \sqrt{30}/6\\
\sqrt2/4 & 0 & \sqrt6/3\\
\sqrt{30}/12 & -\sqrt6/3 & 0
\end{pmatrix}.
\label{eq:wp15-A}
\end{equation}
For a POVM effect $M_y$, let
\begin{equation}
p_y=\Tr(\rho_0M_y),\qquad z_y=\Tr(AM_y).
\end{equation}
The two-quadrature classical Fisher trace is
\begin{equation}
\Tr F_1=\sum_{y:p_y>0}\frac{|z_y|^2}{p_y}.
\end{equation}

\paragraph{Rank-one refinement.}
If $M=\sum_jM_j$ is decomposed into positive effects and $p_j=\Tr(\rho_0M_j)$, $z_j=\Tr(AM_j)$, then quadratic-over-linear convexity gives
\begin{equation}
\frac{|\sum_jz_j|^2}{\sum_jp_j}
\le\sum_j\frac{|z_j|^2}{p_j}.
\end{equation}
Therefore refining effects into rank-one components cannot decrease the Fisher trace. It is sufficient to establish a pointwise bound for $M=w|\phi\rangle\langle\phi|$.

\paragraph{Dual quadratic witness.}
Define
\begin{equation}
Y=
\begin{pmatrix}
9/16 & 0 & 0\\
0 & 9/16 & 3\sqrt{15}/8\\
0 & 3\sqrt{15}/8 & 23/4
\end{pmatrix}.
\label{eq:wp15-Y}
\end{equation}
We claim
\begin{equation}
|\langle\phi|A|\phi\rangle|^2
\le
\langle\phi|\rho_0|\phi\rangle
\langle\phi|Y|\phi\rangle
\label{eq:wp15-witness}
\end{equation}
for every vector $|\phi\rangle$.

For $\lambda>0$ and phase $\theta$, define
\begin{equation}
\mathcal M(\lambda,\theta)
=
\lambda\rho_0+\lambda^{-1}Y
-\left(e^{i\theta}A+e^{-i\theta}A^\dagger\right).
\label{eq:wp15-LMI}
\end{equation}
Set $x=\lambda^2$ and $t=\cos^2\theta$. Direct symbolic reduction gives the first leading principal minor
\begin{equation}
\frac{8x+9}{16\lambda}>0,
\end{equation}
the second leading principal minor
\begin{equation}
\frac{
 t(8x-9)^2+(1-t)(64x^2+112x+81)
}{256x},
\label{eq:wp15-minor2}
\end{equation}
and the determinant
\begin{equation}
\frac{
 t(8x-9)^2+(1-t)(40x+81)
}{128\lambda^3}.
\label{eq:wp15-det}
\end{equation}
For $x>0$ and $0\le t\le1$, both numerators are manifestly nonnegative. They are strictly positive except at the isolated point $t=1$, $x=9/8$; positive semidefiniteness there follows by continuity. Hence
\begin{equation}
\mathcal M(\lambda,\theta)\succeq0
\qquad\text{for all }\lambda>0,\ \theta\in\mathbb R.
\end{equation}
For an arbitrary $|\phi\rangle$, choose $\theta$ to align the phase of $\langle\phi|A|\phi\rangle$. Then
\begin{equation}
2|\langle\phi|A|\phi\rangle|
\le
\lambda\langle\phi|\rho_0|\phi\rangle
+\lambda^{-1}\langle\phi|Y|\phi\rangle.
\end{equation}
Minimizing over $\lambda$ proves Eq.~\eqref{eq:wp15-witness}.

For a rank-one POVM, Eq.~\eqref{eq:wp15-witness} gives
\begin{equation}
\Tr F_1
\le
\sum_y w_y\langle\phi_y|Y|\phi_y\rangle
=\Tr Y
=\frac{55}{8}.
\label{eq:wp15-upper}
\end{equation}
Rank-one refinement extends the upper bound to every POVM.

\paragraph{Sharp limiting projective sequence.}
Write Eq.~\eqref{eq:wp15-A} in block form
\begin{equation}
A=\begin{pmatrix}B&b\\a^\dagger&0\end{pmatrix},
\end{equation}
where
\begin{equation}
B=
\begin{pmatrix}
0&\sqrt2/2\\
\sqrt2/4&0
\end{pmatrix},
\quad
 a=\left(\frac{\sqrt{30}}{12},-\frac{\sqrt6}{3}\right),
\quad
 b=\begin{pmatrix}\sqrt{30}/6\\\sqrt6/3\end{pmatrix}.
\end{equation}
The numerical radius of $B$ is
\begin{equation}
w(B)=\frac{3\sqrt2}{8}.
\end{equation}
The two orthogonal support vectors
\begin{equation}
|s_\pm\rangle=\frac{|e_1\rangle\pm|e_2\rangle}{\sqrt2}
\end{equation}
therefore contribute a total
\begin{equation}
\frac98
\end{equation}
to the Fisher trace.

For a vector approaching the kernel,
\begin{equation}
|\phi_\epsilon\rangle
=\cos\epsilon|q\rangle+\sin\epsilon|t\rangle,
\end{equation}
the baseline probability is $\sin^2\epsilon/2$, while its Fisher contribution tends to
\begin{equation}
2|\langle q|A|t\rangle+\langle t|A|q\rangle|^2.
\end{equation}
Choosing
\begin{equation}
|t_*\rangle=i\frac{(a-b)^T}{\|a-b\|}
\end{equation}
gives
\begin{equation}
\|a-b\|^2=\frac{23}{8},
\end{equation}
so the near-kernel contribution tends to
\begin{equation}
\frac{23}{4}.
\end{equation}
Let $U_\epsilon$ be a unitary rotation by angle $\epsilon$ in $\operatorname{span}\{|q\rangle,|t_*\rangle\}$ and the identity on its orthogonal complement. The projective measurement
\begin{equation}
\left\{
U_\epsilon|s_+\rangle,
U_\epsilon|s_-\rangle,
U_\epsilon|q\rangle
\right\}
\end{equation}
has ordinary positive baseline probabilities for every sufficiently small $\epsilon>0$, and
\begin{equation}
\lim_{\epsilon\to0^+}\Tr F_1[U_\epsilon]
=\frac98+\frac{23}{4}
=\frac{55}{8}.
\label{eq:wp15-lower}
\end{equation}
Combining Eqs.~\eqref{eq:wp15-upper} and \eqref{eq:wp15-lower},
\begin{equation}
\boxed{
\sup_{\mathrm{POVM}}\Tr F_1=\frac{55}{8}.
}
\end{equation}
The result is a supremum; exact attainment by a regular POVM is not required.

For this benchmark the three distinct layers are therefore
\begin{equation}
12\quad>\quad\frac{43}{4}\quad>\quad\frac{55}{8},
\end{equation}
corresponding respectively to the physical shared-curvature resource ceiling, the SLD-QFI trace, and the exact accessible one-copy common-record Fisher supremum.
